\section{Definiciones previas}

\subsection{Pruebas post-hoc}
Sea $X$ un conjunto de datos aleatorios generado bajo un modelo estadístico $\{P_\theta:\theta\in\Theta\}$. 
Considérese una \emph{familia} de hipótesis nulas
\[
\mathcal{H}=\{H_{0,1},H_{0,2},\dots,H_{0,m}\},
\]
donde cada $H_{0,i}$ impone una restricción sobre $\theta$ (por ejemplo, igualdad de parámetros entre grupos o independencia entre variables categóricas). 

Un procedimiento de prueba múltiple puede representarse mediante un vector de decisiones
\[
\boldsymbol{\delta}(X)=\big(\delta_1(X),\dots,\delta_m(X)\big), 
\qquad \delta_i(X)\in\{0,1\},
\]
donde $\delta_i(X)=1$ indica rechazo de $H_{0,i}$ y $\delta_i(X)=0$ indica no rechazo.

Se denominan \textbf{pruebas post-hoc} a los contrastes aplicados a una familia $\mathcal{H}$ \emph{después} de una etapa previa basada en los datos (por ejemplo, un contraste global o un análisis exploratorio) que motiva investigar \emph{hipótesis específicas} dentro de $\mathcal{H}$. 


\subsection{Falso positivo}
Considérese una única hipótesis nula $H_0$ con espacio paramétrico nulo $\Theta_0\subseteq\Theta$ y una regla de decisión $\delta(X)\in\{0,1\}$, donde $\delta(X)=1$ denota el rechazo de $H_0$.

Un \textbf{falso positivo} ocurre cuando el procedimiento rechaza una hipótesis nula que es verdadera. Formalmente, el evento de falso positivo es
\[
\{\delta(X)=1\}\cap\{\theta\in\Theta_0\}.
\]
La probabilidad de falso positivo (error tipo I) bajo el parámetro $\theta$ es
\[
P_\theta\big(\delta(X)=1\big), \qquad \theta\in\Theta_0.
\]

El \textbf{tamaño} (o \emph{size}) de una prueba es la \emph{máxima} probabilidad de cometer un falso positivo cuando la hipótesis nula es verdadera. Para definirlo formalmente, sea $\Theta_0$ el conjunto de valores del parámetro $\theta$ para los cuales $H_0$ es cierta. Entonces, para cada $\theta\in\Theta_0$ existe una probabilidad de rechazar $H_0$ aun siendo verdadera:
\[
P_\theta\big(\delta(X)=1\big).
\]
Como dicha probabilidad puede cambiar según el valor específico de $\theta$ dentro de $\Theta_0$, se toma el peor caso (el más grande) sobre todos los escenarios compatibles con $H_0$. A esa máxima probabilidad se le llama \textbf{tamaño}:
\[
\alpha \;=\;\sup_{\theta\in\Theta_0} P_\theta\big(\delta(X)=1\big).
\]

\textbf{Interpretación de $\sup$ (supremo).} El símbolo $\sup$ significa \emph{supremo} o \emph{cota superior mínima}. En este contexto, puede leerse como:
\[
\sup_{\theta\in\Theta_0} P_\theta(\delta(X)=1) \;=\; \text{``la mayor probabilidad posible de rechazar $H_0$ cuando $H_0$ es verdadera''}.
\]
Si esa mayor probabilidad se alcanza en algún valor $\theta^\star\in\Theta_0$, entonces el supremo coincide con un máximo:
\[
\sup_{\theta\in\Theta_0} P_\theta(\delta(X)=1)=\max_{\theta\in\Theta_0} P_\theta(\delta(X)=1).
\]
Si no se alcanza exactamente, el supremo sigue representando el límite superior más ajustado (la mejor cota superior).

\textbf{Interpretación de $\alpha$.} Decir que una prueba es de \emph{nivel} $\alpha$ significa que, bajo cualquier situación en la que $H_0$ sea verdadera (es decir, para todo $\theta\in\Theta_0$),
\[
P_\theta\big(\delta(X)=1\big)\le \alpha.
\]
En palabras: \emph{aunque el valor real de $\theta$ bajo $H_0$ sea el peor posible, la probabilidad de falso positivo no excede $\alpha$.}
Por ejemplo, si $\alpha=0.05$, la prueba está construida para que la probabilidad de rechazar $H_0$ cuando es cierta sea, en el peor caso, como máximo $5\%$.

\subsection{Family-Wise Error Rate (FWER)}
Sea $\mathcal{H}=\{H_{0,1},\dots,H_{0,m}\}$ una familia de hipótesis y $\boldsymbol{\delta}(X)=(\delta_1(X),\dots,\delta_m(X))$ el vector de decisiones asociado. 
Defínase el conjunto de índices de hipótesis nulas verdaderas como
\[
I_0 \subseteq \{1,\dots,m\}.
\]
El evento ``cometer al menos un falso positivo dentro de la familia'' es
\[
\bigcup_{i\in I_0}\{\delta_i(X)=1\}.
\]
El \textbf{Family-Wise Error Rate (FWER)} se define como
\[
\mathrm{FWER}
\;=\;
P\!\left(\bigcup_{i\in I_0}\{\delta_i(X)=1\}\right).
\]
En términos conceptuales, $\mathrm{FWER}$ es la probabilidad de realizar \emph{al menos un} rechazo incorrecto (error tipo I) dentro del conjunto de comparaciones.
